% =====================================================================
\chapter{El modo directo y los números duales}
\label{tema:2}
% =====================================================================

El modo directo propaga, junto con cada valor de la traza, su derivada respecto a una dirección fijada de las entradas. Este tema lo formula como propagación de tangentes, demuestra que calcula un producto jacobiano-vector y lo reinterpreta con los números duales, que reducen la regla de la cadena a aritmética. Se siguen la §3.1 de~\cite{baydin}, incluida la §3.1.1 sobre números duales, el capítulo~3 de~\cite{gw}, las notas de la clase~8 de MIT 18.337~\cite{mit18337} y la clase~5 de MIT 18.S096~\cite{mit18s096}.

\section{Propagación de tangentes}

Se fija una traza de $f\colon\R^{n}\to\R^{m}$ (\cref{def:traza}), un punto $x$ y una dirección $\dot x\in\R^{n}$.

\begin{definicion}[Tangentes]\label{def:tangente}
Las \emph{tangentes} de la traza son los números definidos por recurrencia como
\[
\dot v_{i-n}=\dot x_i\quad(i=1,\dots,n),\qquad
\dot v_i=\sum_{j\pre i}\frac{\partial\varphi_i}{\partial v_j}\,\dot v_j\quad(i=1,\dots,\ell),
\]
donde las derivadas parciales se evalúan en los valores de la traza. El vector $\dot x$ se llama \emph{semilla}.
\end{definicion}

\begin{teorema}[El modo directo calcula $J\dot x$]\label{teo:directo}
Sea $\gamma(t)=x+t\dot x$. Para cada $i$, $\dot v_i=\frac{d}{dt}v_i(\gamma(t))\big|_{t=0}=\nabla v_i(x)\T\dot x$, donde $v_i(\cdot)$ denota la variable $i$ como función de las entradas. En particular, el vector de tangentes de las salidas es $\dot y=J_f(x)\,\dot x$.
\end{teorema}

\begin{proof}
Por inducción sobre $i$. Para las entradas, $v_{i-n}(\gamma(t))=x_i+t\dot x_i$, cuya derivada en $t=0$ es $\dot x_i$. Sea $i\ge 1$ y supóngase el resultado cierto para todo índice menor; en particular, para cada $j\pre i$, porque $j<i$. Como $v_i(\gamma(t))=\varphi_i\bigl((v_j(\gamma(t)))_{j\pre i}\bigr)$, la regla de la cadena (módulo~4) da
\[
\frac{d}{dt}v_i(\gamma(t))\Big|_{t=0}=\sum_{j\pre i}\frac{\partial\varphi_i}{\partial v_j}\,\frac{d}{dt}v_j(\gamma(t))\Big|_{t=0}=\sum_{j\pre i}\frac{\partial\varphi_i}{\partial v_j}\,\dot v_j=\dot v_i.
\]
Para las salidas, $\frac{d}{dt}f(x+t\dot x)|_{t=0}=J_f(x)\dot x$ es la derivada direccional.
\end{proof}

\begin{ejemplo}[Traza tangente]\label{ej:tangente}
En el \cref{ej:baydin}, con semilla $\dot x=(1,0)\T$ se obtiene $\partial y/\partial x_1$~\cite[tabla~2]{baydin}:
\begin{center}\small
\begin{tabular}{@{}lll@{}}
\toprule
Tangente & Regla & Valor \\
\midrule
$\dot v_{-1}$, $\dot v_0$ & semilla & $1$, $0$ \\
$\dot v_1$ & $\dot v_{-1}/v_{-1}$ & $0{,}5$ \\
$\dot v_2$ & $\dot v_{-1}v_0+v_{-1}\dot v_0$ & $5$ \\
$\dot v_3$ & $\cos(v_0)\,\dot v_0$ & $0$ \\
$\dot v_4$ & $\dot v_1+\dot v_2$ & $5{,}5$ \\
$\dot v_5$ & $\dot v_4-\dot v_3$ & $5{,}5$ \\
\bottomrule
\end{tabular}
\end{center}
En efecto, $\partial f/\partial x_1=1/x_1+x_2=0{,}5+5=5{,}5$. Con la semilla $(0,1)\T$ se obtendría $\partial f/\partial x_2=x_1-\cos x_2=2-\cos 5\approx 1{,}7163$.
\end{ejemplo}

\begin{corolario}[Coste del jacobiano en modo directo]\label{cor:coste-directo}
Una pasada en modo directo da una columna de $J_f(x)$, la correspondiente a la semilla $e_k$. El jacobiano completo de $f\colon\R^{n}\to\R^{m}$ requiere $n$ pasadas, independientemente de $m$.
\end{corolario}

\begin{proof}
Por el \cref{teo:directo}, con $\dot x=e_k$ se obtiene $J_f(x)e_k$, que es la columna $k$. Las $n$ semillas $e_1,\dots,e_n$ dan las $n$ columnas.
\end{proof}

El corolario fija el campo de aplicación del modo directo: es óptimo cuando hay pocas entradas y muchas salidas y es inviable para el gradiente de una pérdida escalar respecto a $10^{11}$ parámetros, que exigiría $10^{11}$ pasadas.

\section{Números duales}

\begin{definicion}[Números duales]\label{def:duales}
El conjunto de los \emph{números duales} es $\Du=\{a+b\varepsilon : a,b\in\R\}$, con la suma y el producto
\[
(a+b\varepsilon)+(c+d\varepsilon)=(a+c)+(b+d)\varepsilon,\qquad
(a+b\varepsilon)(c+d\varepsilon)=ac+(ad+bc)\varepsilon .
\]
El número $a$ es la \emph{parte real} y $b$, la \emph{parte dual}.
\end{definicion}

El producto es el que resulta de multiplicar formalmente e imponer $\varepsilon^{2}=0$ con $\varepsilon\ne 0$. Con estas operaciones $\Du$ es un anillo conmutativo con unidad, isomorfo a $\R[\varepsilon]/(\varepsilon^{2})$, pero no es un cuerpo: $\varepsilon\cdot\varepsilon=0$ con $\varepsilon\ne 0$, de modo que $\varepsilon$ no tiene inverso~\cite[§3.1.1]{baydin}.

\begin{proposicion}[Inverso de un dual]\label{prop:inverso}
Si $a\ne 0$, el dual $a+b\varepsilon$ es invertible y
\[
\frac{1}{a+b\varepsilon}=\frac1a-\frac{b}{a^{2}}\,\varepsilon .
\]
\end{proposicion}

\begin{proof}
$(a+b\varepsilon)\bigl(\tfrac1a-\tfrac{b}{a^{2}}\varepsilon\bigr)=1+\bigl(-\tfrac{b}{a}+\tfrac{b}{a}\bigr)\varepsilon=1$. Si $a=0$, $(b\varepsilon)(c+d\varepsilon)=bc\,\varepsilon\ne 1$ para todo $c,d$, y no hay inverso.
\end{proof}

La parte dual del inverso es $b$ veces la derivada de $u\mapsto 1/u$ en $a$, y la parte dual del producto es la regla de Leibniz. Esto no es casual, y motiva la definición siguiente.

\begin{definicion}[Extensión dual de una operación elemental]\label{def:extension}
Sea $\varphi\colon U\subseteq\R^{k}\to\R$ de clase $C^{1}$. Su \emph{extensión dual} es
\[
\varphi(a_1+b_1\varepsilon,\dots,a_k+b_k\varepsilon)=\varphi(a)+\Bigl(\sum_{j=1}^{k}\frac{\partial\varphi}{\partial u_j}(a)\,b_j\Bigr)\varepsilon ,\qquad a\in U.
\]
Por ejemplo, $\exp(a+b\varepsilon)=e^{a}+e^{a}b\,\varepsilon$ y $\sen(a+b\varepsilon)=\sen a+(\cos a)\,b\,\varepsilon$.
\end{definicion}

\begin{proposicion}[Coherencia con la aritmética]\label{prop:coherencia}
Para la suma, el producto y la inversión, la extensión dual de la \cref{def:extension} coincide con las operaciones de la \cref{def:duales} y de la \cref{prop:inverso}.
\end{proposicion}

\begin{proof}
Para $\varphi(u_1,u_2)=u_1+u_2$, las parciales valen~1 y la extensión da $(a_1+a_2)+(b_1+b_2)\varepsilon$. Para $\varphi(u_1,u_2)=u_1u_2$, las parciales son $a_2$ y $a_1$, y la extensión da $a_1a_2+(a_2b_1+a_1b_2)\varepsilon$, que es el producto de la \cref{def:duales}. Para $\varphi(u)=1/u$, $\varphi'(a)=-1/a^{2}$ y la extensión da la fórmula de la \cref{prop:inverso}.
\end{proof}

\begin{teorema}[Los duales propagan la derivada]\label{teo:duales}
Sea $f\colon\R^{n}\to\R^{m}$ dada por una traza y evalúese la traza en aritmética dual, con entradas $x_k+\dot x_k\varepsilon$ y la extensión dual de cada operación elemental. Entonces cada variable vale $v_i+\dot v_i\varepsilon$, con $\dot v_i$ la tangente de la \cref{def:tangente}, y las salidas valen
\[
f(x+\dot x\,\varepsilon)=f(x)+J_f(x)\,\dot x\,\varepsilon .
\]
\end{teorema}

\begin{proof}
Por inducción sobre $i$. Para las entradas es la elección de la semilla. Si las variables $v_j$ con $j\pre i$ valen $v_j+\dot v_j\varepsilon$, la \cref{def:extension} da
$\varphi_i\bigl((v_j+\dot v_j\varepsilon)_{j\pre i}\bigr)=\varphi_i\bigl((v_j)_{j\pre i}\bigr)+\bigl(\sum_{j\pre i}\tfrac{\partial\varphi_i}{\partial v_j}\dot v_j\bigr)\varepsilon=v_i+\dot v_i\varepsilon$. La última afirmación es el \cref{teo:directo}.
\end{proof}

El teorema justifica la lectura del desarrollo de Taylor del módulo~4 con un infinitesimal: $f(a+b\varepsilon)=f(a)+f'(a)b\varepsilon$ es el desarrollo de primer orden, y aquí el truncamiento es exacto porque los términos de orden superior contienen $\varepsilon^{2}=0$. Para evaluar $f'(a)$ basta con calcular $f(a+\varepsilon)$ y leer la parte dual.

\begin{ejemplo}\label{ej:dual-cubo}
Sea $f(x)=x^{3}+2x$ en $x=2$. Con $x=2+\varepsilon$:
\begin{enumerate}[label=\arabic*.]
  \item $x^{2}=(2+\varepsilon)(2+\varepsilon)=4+4\varepsilon$;
  \item $x^{3}=(4+4\varepsilon)(2+\varepsilon)=8+(4+8)\varepsilon=8+12\varepsilon$;
  \item $2x=4+2\varepsilon$, y $f=12+14\varepsilon$.
\end{enumerate}
Luego $f(2)=12$ y $f'(2)=14$, que coincide con $f'(x)=3x^{2}+2$.
\end{ejemplo}

\begin{observacion}[Duales con varias componentes]\label{obs:multidual}
Si se toman $p$ infinitesimales $\varepsilon_1,\dots,\varepsilon_p$ con $\varepsilon_r\varepsilon_s=0$ para todo $r,s$, cada variable lleva un vector de $p$ partes duales, y con las semillas $x_k+\sum_r(\dot x^{(r)})_k\varepsilon_r$ una sola pasada calcula $p$ productos $J\dot x^{(r)}$. Con $p=n$ y semillas $e_1,\dots,e_n$ se obtiene el jacobiano entero en una pasada vectorial, cuyo coste sigue siendo proporcional a $n\cdot\ops(f)$~\cite{mit18337}. La biblioteca JAX ofrece esta operación como \texttt{jvp} (\emph{Jacobian-vector product}).
\end{observacion}

\section{Coste del modo directo}

\begin{proposicion}\label{prop:coste-dual}
Existe una constante $\omega_{\mathrm{dir}}$, que depende solo del conjunto $\Phi$ de operaciones elementales, tal que una pasada en modo directo cuesta a lo sumo $\omega_{\mathrm{dir}}\cdot\ops(f)$ operaciones.
\end{proposicion}

\begin{proof}
Cada operación de la traza se sustituye por su extensión dual. La parte real cuesta lo mismo que la operación original; la parte dual requiere evaluar a lo sumo $k$ derivadas parciales, de coste acotado por hipótesis (\cref{def:elementales}), y una combinación lineal de $k$ términos. El cociente entre ambos costes está acotado por una constante que depende de $k$ y de $\Phi$, y sumando sobre las $\ell$ operaciones se obtiene la cota. Por ejemplo, el producto dual usa tres productos y una suma en lugar de un producto.
\end{proof}

\begin{observacion}[Conexión con los LLM]
El modo directo no se usa para el gradiente de la pérdida de un LLM, pero sí aparece en otros cálculos: un producto hessiana-vector se obtiene aplicando el modo directo al gradiente (\cref{tema:5}), y la memoria del modo directo no crece con la longitud de la traza, porque cada tangente se descarta en cuanto deja de usarse. El precio es el del \cref{cor:coste-directo}: una pasada por dirección.
\end{observacion}

\begin{lecturas}
\item Baydin \emph{et al.}~\cite{baydin}, §3.1 «Forward Mode» y §3.1.1 «Dual Numbers».
\item Griewank y Walther~\cite{gw}, cap.~3 «Fundamentals of Forward and Reverse».
\item MIT 18.337~\cite{mit18337}, notas de la clase~8, apartados «Dual numbers», «Higher dimensions» y «Forward-Mode AD as jvp».
\item MIT 18.S096~\cite{mit18s096}, clase~5, parte~2, «Forward Automatic Differentiation via Dual Numbers».
\end{lecturas}

\section*{Ejercicios del tema 2}
\addcontentsline{toc}{section}{Ejercicios del tema 2}

\begin{ejercicio}{1}{eje:2-cuadrado}
(a) Calcule la parte dual de $(3+\varepsilon)^{2}$ e interprétela. (b) Una función $f\colon\R^{1000}\to\R$ se diferencia en modo directo. ¿Cuántas pasadas hacen falta para obtener su gradiente? ¿Y para una sola derivada direccional?
\end{ejercicio}

\begin{ejercicio}{2}{eje:2-cuartica}
Con números duales, calcule $f(1)$ y $f'(1)$ para $f(x)=x^{4}-3x$.
\end{ejercicio}

\begin{ejercicio}{2}{eje:2-cociente}
Con números duales y la \cref{prop:inverso}, calcule $g(2)$ y $g'(2)$ para $g(x)=x/(1+x^{2})$.
\end{ejercicio}

\begin{ejercicio}{2}{eje:2-tanh}
Con la extensión dual de $\tanh$, calcule $h(1)$ y $h'(1)$ para $h(x)=\tanh(x^{2})$ con cuatro cifras significativas, y compruébelo con la regla de la cadena.
\end{ejercicio}

\begin{ejercicio}{2}{eje:2-multidual}
Para $f(x_1,x_2)=\ln x_1+x_1x_2-\sen x_2$ en $(2,5)$, evalúe la traza del \cref{ej:baydin} con duales de dos componentes, $x_1=2+\varepsilon_1$ y $x_2=5+\varepsilon_2$ ($\varepsilon_r\varepsilon_s=0$), y obtenga el gradiente en una sola pasada. Deduzca la derivada direccional en la dirección $(1,1)\T$.
\fuente{adaptado de MIT 18.337~\cite{mit18337}, notas de la clase~8, apartado «Higher dimensions», aplicado al ejemplo de Baydin \emph{et al.}~\cite[tabla~2]{baydin}.}
\end{ejercicio}

\begin{ejercicio}{3}{eje:2-polinomios}
Demuestre por inducción sobre el grado, usando solo la aritmética de la \cref{def:duales}, que todo polinomio real $p$ cumple $p(a+b\varepsilon)=p(a)+p'(a)\,b\,\varepsilon$. Demuestre después que, si $f$ y $g$ cumplen esa propiedad, también la cumple $g\circ f$.
\end{ejercicio}
